% --- Archon physics formalization source begin ---
% archon:physics
% archon:covers PhyXMiniProblems/problem_phyx_mini_0060.lean
% archon:source-report reports/phyx_mini/problem_phyx_mini_0060.source.json

\chapter{Physics problem phyx\_mini\_0060}
\label{ch:PhyXMiniProblems_problem_phyx_mini_0060}

\paragraph{Problem source.}
\#\# Physical scenario

The mirror in figure deflects a horizontal laser beam by 60\textasciicircum{}\textbackslash{}circ.

\#\# Question

What is the angle \textbackslash{}phi?

\#\# Answer choices

- A: A: 30\textasciicircum{}\textbackslash{}circ

- B: B: 90\textasciicircum{}\textbackslash{}circ

- C: C: 60\textasciicircum{}\textbackslash{}circ

- D: D: 120\textasciicircum{}\textbackslash{}circ

\#\# Figure caption (auxiliary; use the image as primary evidence)

The image is a geometric figure featuring angles and lines. There's a horizontal purple line intersected by another line at an angle, forming a right angle and a supplementary angle. The intersecting line is depicted with an arrow pointing upwards. The angle formed between the lines is labeled as \textbackslash{}(60\textasciicircum{}\textbackslash{}circ\textbackslash{}).

Additionally, there is a shaded slanted surface beneath where the lines intersect, forming an angle labeled as \textbackslash{}(\textbackslash{}phi\textbackslash{}) with the horizontal line. This configuration typically represents a situation involving forces or vectors in physics, where \textbackslash{}(\textbackslash{}phi\textbackslash{}) might denote an angle of inclination or a related angle in a problem scenario.

\#\# Dataset metadata

Geometrical Optics; Spatial Relation Reasoning

\paragraph{Recorded answer/context.}
C: C: 60\textasciicircum{}\textbackslash{}circ

\paragraph{Figure/image path.}
/root/proposal\_for\_physic/science-mango/phyx\_mini\_run/phyx\_data/test\_image/60.png

\paragraph{Formalization target.}
create a compiling Lean file with sorry bodies at `PhyXMiniProblems/problem\_phyx\_mini\_0060.lean`. The Lean declarations must preserve the physical quantities, dimensions or dimensional roles, figure labels, governing-law hypotheses, and final relation expressed by this problem.
Use Mathlib/Physlib names found through LeanExplore where available. If a physics API is missing, introduce faithful local abstractions rather than scalar placeholder aliases.

\begin{theorem}[Physics formalization target]
\label{thm:physics:phyx_mini_0060:target}
\lean{PhyXMiniProblems.ProblemPhyXMini0060.problem_phyx_mini_0060}
\uses{def:physics:phyx-mini-0060:phyxminiproblems-problemphyxmini0060-degrees, def:physics:phyx-mini-0060:phyxminiproblems-problemphyxmini0060-mirrordeflectionsetup, def:physics:phyx-mini-0060:phyxminiproblems-problemphyxmini0060-hasgeometricanglereadouts, def:physics:phyx-mini-0060:phyxminiproblems-problemphyxmini0060-matchesfigurereadouts, def:physics:phyx-mini-0060:phyxminiproblems-problemphyxmini0060-obeysspecularreflectionlaw, def:physics:phyx-mini-0060:phyxminiproblems-problemphyxmini0060-matchesanswer}
The pictured mirror turns a horizontal incident ray through \(60^\circ\).
Specular reflection therefore places the mirror tangent at
\(\phi=30^\circ\), answer A; this corrects the dataset's inconsistent recorded
choice C.
\end{theorem}
\begin{proof}
The figure gives incident direction angle \(0\) and reflected direction angle
\(60^\circ\).  On the acute branch, the angle form of the reflection law says
\[
 60^\circ=2\phi-0.
\]
Thus \(\phi=30^\circ\).  Expanding the answer map shows that choice A is
exactly the same radian value.
\end{proof}

% --- Archon declaration topology begin ---
\section{Declaration topology}
\begin{definition}[Plane]
  \label{def:physics:phyx-mini-0060:phyxminiproblems-problemphyxmini0060-plane}
  \lean{PhyXMiniProblems.ProblemPhyXMini0060.Plane}
  The Euclidean plane containing the mirror and both laser rays
\end{definition}
\begin{proof}
  This is the declared model definition of Plane.
\end{proof}
\begin{definition}[degrees]
  \label{def:physics:phyx-mini-0060:phyxminiproblems-problemphyxmini0060-degrees}
  \lean{PhyXMiniProblems.ProblemPhyXMini0060.degrees}
  Convert a dimensionless angle readout in degrees to radians
\end{definition}
\begin{proof}
  This is the declared model definition of degrees.
\end{proof}
\begin{definition}[Light Ray]
  \label{def:physics:phyx-mini-0060:phyxminiproblems-problemphyxmini0060-lightray}
  \lean{PhyXMiniProblems.ProblemPhyXMini0060.LightRay}
  \uses{def:physics:phyx-mini-0060:phyxminiproblems-problemphyxmini0060-plane}
  An oriented laser ray. Its direction is normalized, so the scalar parameter of pointAt has the same length readout as the plane coordinates
\end{definition}
\begin{proof}
  This is the declared model definition of Light Ray.
\end{proof}
\begin{definition}[point At]
  \label{def:physics:phyx-mini-0060:phyxminiproblems-problemphyxmini0060-lightray-pointat}
  \lean{PhyXMiniProblems.ProblemPhyXMini0060.LightRay.pointAt}
  \uses{def:physics:phyx-mini-0060:phyxminiproblems-problemphyxmini0060-plane, def:physics:phyx-mini-0060:phyxminiproblems-problemphyxmini0060-lightray}
  The point at the given signed distance along a laser ray
\end{definition}
\begin{proof}
  This is the declared model definition of point At.
\end{proof}
\begin{definition}[Plane Mirror]
  \label{def:physics:phyx-mini-0060:phyxminiproblems-problemphyxmini0060-planemirror}
  \lean{PhyXMiniProblems.ProblemPhyXMini0060.PlaneMirror}
  \uses{def:physics:phyx-mini-0060:phyxminiproblems-problemphyxmini0060-plane}
  An ideal plane mirror in the two-dimensional optical diagram. The named unit tangent points toward the right in the figure and fixes the acute representative of the mirror's otherwise unoriented surface line
\end{definition}
\begin{proof}
  This is the declared model definition of Plane Mirror.
\end{proof}
\begin{definition}[reflect Direction]
  \label{def:physics:phyx-mini-0060:phyxminiproblems-problemphyxmini0060-planemirror-reflectdirection}
  \lean{PhyXMiniProblems.ProblemPhyXMini0060.PlaneMirror.reflectDirection}
  \uses{def:physics:phyx-mini-0060:phyxminiproblems-problemphyxmini0060-plane, def:physics:phyx-mini-0060:phyxminiproblems-problemphyxmini0060-planemirror}
  Reflect a propagation direction in the affine mirror. Translating the direction to the impact point allows direct use of Mathlib's affine reflection
\end{definition}
\begin{proof}
  This is the declared model definition of reflect Direction.
\end{proof}
\begin{definition}[Mirror Deflection Setup]
  \label{def:physics:phyx-mini-0060:phyxminiproblems-problemphyxmini0060-mirrordeflectionsetup}
  \lean{PhyXMiniProblems.ProblemPhyXMini0060.MirrorDeflectionSetup}
  \uses{def:physics:phyx-mini-0060:phyxminiproblems-problemphyxmini0060-plane, def:physics:phyx-mini-0060:phyxminiproblems-problemphyxmini0060-lightray, def:physics:phyx-mini-0060:phyxminiproblems-problemphyxmini0060-planemirror}
  The physical objects and labeled angular readouts in the source figure. phiRadians is the angle labeled φ between the horizontal reference and the mirror surface. It is an angle readout, not a scalar replacement for the mirror or either laser ray
\end{definition}
\begin{proof}
  This is the declared model definition of Mirror Deflection Setup.
\end{proof}
\begin{definition}[Has Geometric Angle Readouts]
  \label{def:physics:phyx-mini-0060:phyxminiproblems-problemphyxmini0060-hasgeometricanglereadouts}
  \lean{PhyXMiniProblems.ProblemPhyXMini0060.HasGeometricAngleReadouts}
  \uses{def:physics:phyx-mini-0060:phyxminiproblems-problemphyxmini0060-mirrordeflectionsetup}
  The three scalar angle labels agree with Mathlib's undirected angles between the corresponding nonzero geometric directions
\end{definition}
\begin{proof}
  This is the declared model definition of Has Geometric Angle Readouts.
\end{proof}
\begin{definition}[Matches Figure Readouts]
  \label{def:physics:phyx-mini-0060:phyxminiproblems-problemphyxmini0060-matchesfigurereadouts}
  \lean{PhyXMiniProblems.ProblemPhyXMini0060.MatchesFigureReadouts}
  \uses{def:physics:phyx-mini-0060:phyxminiproblems-problemphyxmini0060-degrees, def:physics:phyx-mini-0060:phyxminiproblems-problemphyxmini0060-lightray-pointat, def:physics:phyx-mini-0060:phyxminiproblems-problemphyxmini0060-mirrordeflectionsetup}
  Figure-derived data: the incoming beam is horizontal and reaches the mirror, the outgoing beam begins at the impact point, and the marked deflection is 60°. Only the acute range of φ is read from the drawing; its requested numerical value is intentionally absent
\end{definition}
\begin{proof}
  This is the declared model definition of Matches Figure Readouts.
\end{proof}
\begin{definition}[Obeys Specular Reflection Law]
  \label{def:physics:phyx-mini-0060:phyxminiproblems-problemphyxmini0060-obeysspecularreflectionlaw}
  \lean{PhyXMiniProblems.ProblemPhyXMini0060.ObeysSpecularReflectionLaw}
  \uses{def:physics:phyx-mini-0060:phyxminiproblems-problemphyxmini0060-planemirror-reflectdirection, def:physics:phyx-mini-0060:phyxminiproblems-problemphyxmini0060-mirrordeflectionsetup}
  The governing law of specular reflection. The first conjunct is the geometric reflection of the propagation vector in the mirror. The second is its angle-readout form on the pictured branch: the mirror direction bisects the incident and reflected direction angles
\end{definition}
\begin{proof}
  This is the declared model definition of Obeys Specular Reflection Law.
\end{proof}
\begin{definition}[Answer Choice]
  \label{def:physics:phyx-mini-0060:phyxminiproblems-problemphyxmini0060-answerchoice}
  \lean{PhyXMiniProblems.ProblemPhyXMini0060.AnswerChoice}
  Labels of the four answers displayed with the problem
\end{definition}
\begin{proof}
  This is the declared model definition of Answer Choice.
\end{proof}
\begin{definition}[answer Angle Radians]
  \label{def:physics:phyx-mini-0060:phyxminiproblems-problemphyxmini0060-answerangleradians}
  \lean{PhyXMiniProblems.ProblemPhyXMini0060.answerAngleRadians}
  \uses{def:physics:phyx-mini-0060:phyxminiproblems-problemphyxmini0060-degrees, def:physics:phyx-mini-0060:phyxminiproblems-problemphyxmini0060-answerchoice}
  The radian value printed beside each answer choice
\end{definition}
\begin{proof}
  This is the declared model definition of answer Angle Radians.
\end{proof}
\begin{definition}[Matches Answer]
  \label{def:physics:phyx-mini-0060:phyxminiproblems-problemphyxmini0060-matchesanswer}
  \lean{PhyXMiniProblems.ProblemPhyXMini0060.MatchesAnswer}
  \uses{def:physics:phyx-mini-0060:phyxminiproblems-problemphyxmini0060-answerchoice, def:physics:phyx-mini-0060:phyxminiproblems-problemphyxmini0060-answerangleradians}
  A mirror-angle readout agrees exactly with the selected answer
\end{definition}
\begin{proof}
  This is the declared model definition of Matches Answer.
\end{proof}
\begin{definition}[recorded Dataset Answer]
  \label{def:physics:phyx-mini-0060:phyxminiproblems-problemphyxmini0060-recordeddatasetanswer}
  \lean{PhyXMiniProblems.ProblemPhyXMini0060.recordedDatasetAnswer}
  \uses{def:physics:phyx-mini-0060:phyxminiproblems-problemphyxmini0060-answerchoice}
  The answer label recorded in the supplied dataset metadata, contrary to the pictured physics
\end{definition}
\begin{proof}
  This is the declared model definition of recorded Dataset Answer.
\end{proof}
% --- Archon declaration topology end ---
% --- Archon physics formalization source end ---
